\section{Introduction}
\label{sec:intro}

\noindent
This paper is the eleventh in a series on the philosophy of intimacy and the theory of justice, and it takes up the empirical task the prior paper deferred. The prior paper established that the practice of a good relation is the cultivation of a field rather than the construction of a structure, that the matter of this field is the everyday sensory and habitual fabric of a shared life, and that a good field satisfies two conditions, the eudaimonic and the just, of which the just constitutes rather than merely constrains the eudaimonic. It established these as conditions, in the abstract, and it deferred to a subsequent paper the question of how a shared life is in fact cultivated, assessed, and sustained in the more empirical register of habit, household, and the quality of a life lived together. The present paper takes up that question.

The question is how flourishing is known, practised, assessed, and reconceived across the course of a shared life. The course of a shared life is treated here as a cycle, and the cycle as the primary object: a recurrence of cognition, practice, evaluation, reflection, and the renewed cognition that follows, turning across the time of a relation from its beginning through its marriage and its later years. The cycle is not a linear sequence of stages completed once. Its moments interpenetrate and feed back upon one another, and one of them, evaluation, runs through all the others as an unbroken undercurrent rather than occupying a single station. The cycle is, moreover, the condition of the relation's existence. A relation in which value ceases to be generated and returned, in which the parties cease to revise their sense of one another, in which practice halts, is a relation whose cycle has stopped; and a relation whose cycle has stopped does not persist in a steady state but decays. The continued turning of the cycle is the condition of the relation's survival, and the character of its turning, whether it spirals upward, repeats at a loss, or winds down, is the difference between a flourishing life, a life of slow catastrophe, and the extinction of the bond.

The frameworks by which flourishing is known and assessed are treated, in this paper, as resources the cycle draws upon rather than as the primary object. Positivism, evidentialism, normative theory, virtue ethics, and phenomenology each furnish an epistemology and a practice of flourishing, and each is drawn upon at the moments of the cycle to which it is suited; but none governs the cycle as a whole, and the absence of a single governing framework is, as the prior papers established, not a defect but the condition of an object no single framework possesses. The frameworks are the symmetry-broken phases, in the sense of the prior paper, of the cognition and practice of flourishing; the cycle is the movement across them.

The paper carries three commitments that run through its length. The first is an operational concession. The paper is a practice, and a practice requires what is operable and what is assessable; it therefore restricts itself, where it must assess, to value as structured into an assessable proxy, and it does so in the awareness that the structured proxy is not the value that the prior papers held to be transversal across the registers and irreducible to symbolisation. This is the first occasion in the series on which an incomplete value is chosen deliberately, and the choice is the honest discharge of the prior result rather than its abandonment: since the value that is real cannot be cashed without occupying the place of the barred Other, the honest course, where assessment is required, is to concede that what is assessed is a proxy and to bear responsibility for its incompleteness. The second commitment concerns the assessment of justice: a relation's distribution of returns is shown to admit a statistical signature, the skewness of the distribution of value returned relative to the coupling of the parties, by which exploitation within an apparently satisfying cycle becomes visible. The third concerns the cycle as a whole: the cycle accumulates a phase, in the sense of the prior paper's holonomy, and the sign of that phase distinguishes the good cycle from the bad, while the persistence of the cycle distinguishes both from extinction.

The paper is the programmatic and foundational text of the part of the series concerned with the praxis of flourishing, and it is written at the length and the level of detail that a foundation requires, in the expectation that the subsequent papers will draw upon the instruments it develops rather than re-deriving them. Its scope is correspondingly broad: it situates its conception within the philosophy of well-being and the theory of justice, develops the cognitive and dynamical apparatus of the lifecycle, treats each moment of the cycle through the several frameworks, develops the quantitative assessment of justice and surveys the existing instruments of measurement, traces the cross-cultural variation of flourishing and its practice, displays the apparatus on an extended case, and meets the objections its approach invites.

The paper proceeds as follows. It surveys the frameworks and states its method (\xref{sec:background}), and situates its conception within the philosophy of well-being and the theory of justice (\xref{sec:philosophy}). It recapitulates the dynamical-systems vocabulary it will employ (\xref{sec:dynsys}); establishes the cognitive and computational basis of appraisal, in which evaluation is shown to be a continuous and largely unconscious inference before it is an explicit measurement (\xref{sec:appraisal}); and specifies the lifecycle as a real cycle whose persistence conditions the relation's survival (\xref{sec:cycle}). It sets out, in an analytic matrix, the contribution of each framework at each moment in both its epistemological and its practical layer (\xref{sec:matrix}), and then treats the moments of the cycle in turn, cognition (\xref{sec:cognition}), practice (\xref{sec:practice}), evaluation (\xref{sec:evaluation}), reflection (\xref{sec:reflection}), and re-cognition (\xref{sec:recognition}), each as the site at which several frameworks are drawn upon for the task proper to it. It considers the cycle as a whole, in its three regimes of flourishing, catastrophe, and extinction (\xref{sec:regimes}); displays the whole apparatus on an extended longitudinal case (\xref{sec:case}); traces the cross-cultural variation of flourishing and of the practices of marriage and family (\xref{sec:crosscultural}); states its foundational role in the series and the instruments it bequeaths to the papers that follow (\xref{sec:foundational}); meets the objections its approach invites (\xref{sec:objections}); and concludes in plural conclusions (\xref{sec:conclusions}). A formal appendix gathers the notation and the key concepts (\xref{sec:notation}) and develops the spectral, decision-theoretic, and dynamical formalisations the main text employs (\xref{sec:appendix}).
